\cotefile{1}
\scanpage{65}
% archive sheet number: 66
troisième espace de Hilbert, si $u\in L(E,F)$, $v\in L(F,G)$, alors on aura

\[
(\overset{n}{\Lambda}v)(\overset{n}{\Lambda}u)=\overset{n}{\Lambda}(v\circ u).\tag*{(10)}
\]

Notons aussi la formule

\[
\begin{aligned}
&\deleted{(\overset{n}{\Lambda}u)^*=\overset{n}{\Lambda}\unsure{u^*}}\\
&\left\{\begin{aligned}
(u_1\wedge\ldots\wedge u_n)^*&=u_1^*\wedge\ldots\wedge u_n^*\\
(\overset{n}{\Lambda}u)^*&=\overset{n}{\Lambda}(u^*)
\end{aligned}\right.
\end{aligned}\tag*{(11)}
\]

\deleted{La première}

Pour l'établir, on peut p. ex. \deleted{partir \unsure{de la formule de polarisation}} \deleted{se borner au cas où \unsure{les $u$ sont} \unsure{supposant} (car $u_1\wedge\ldots\wedge u_n$ s'exprime \unsure{de façon naturelle} comme combinaison linéaire de \unsure{puissances} $\overset{n}{\Lambda}u$, où $u$ est une somme} partir \deleted{\unsure{du}} (car $u_1\wedge\ldots\wedge u_n=\frac1{n!}\deleted{\sum_{\unsure{p=0}}^{\unsure{n}}(-1)^p}\sum_{\substack{\deleted{\unsure{\sum}(\unsure{i_1,\ldots,}i_p)}\\0\leq p\leq n\\1\leq i_1\unsure{,\ldots,}i_{\unsure{n-p|p}}\leq n}}(-1)^p\deleted{a}\overset{n}{\Lambda}(u_{i_1}+\ldots+u_{i_{n-p}})$ -- voir Introduction, 2, formule (8,--) ), où on a

\[
\begin{aligned}
((\overset{n}{\Lambda}u)\circ a_1\wedge\ldots\wedge a_n,b_1\wedge\ldots\wedge b_n)&=\det((u.a_i,b_j))=\deleted{\unsure{d}}\det((a_i,u^*b_j))\\
&=(a_1\wedge\ldots\wedge a_n,(\overset{n}{\Lambda}u^*).\,b_1\wedge\ldots\wedge b_n),
\end{aligned}
\]

ce qui établit \deleted{que} la deuxième formule (11).

En particulier, supposons $E=F$, alors on voit que l'application $u\longrightarrow\overset{n}{\Lambda}u$ de $L(E,F)$ dans $L(\overset{n}{\Lambda}^{(2)}E,\overset{n}{\Lambda}^{(2)}F)$ est \uline{multiplicative} (formule (10)) \uline{commute avec l'opération $*$} (formule 11) \inserted{elle transforme donc hermitien en hermitien,} \deleted{\unsure{si} De} plus elle transforme \deleted{le} un hermitien \uline{positif} $h$ en un hermitien positif, et on a, pour \deleted{\unsure{tout|$\Lambda$}} un exposant \deleted{\unsure{réel}} $\alpha>0$

\[
(\overset{n}{\Lambda}h)^\alpha=\overset{n}{\Lambda}(h^\alpha)\qquad(h\in L(E),\ h\geq0)\tag*{(12)}
\]

En effet, si $h$ est positif, on a $h=k^2$, où $k$ est un hermitien positif, d'où $(\overset{n}{\Lambda}h)=\overset{n}{\Lambda}k^2=(\overset{n}{\Lambda}k)^2$, où $\overset{n}{\Lambda}k$ est hermitien \deleted{(car \unsure{l'application} \illegible{2}}
